\documentclass[11pt]{article}
\usepackage[a4paper,margin=1in]{geometry}
\usepackage[T1]{fontenc}
\usepackage{lmodern,microtype}
\usepackage{amsmath,amssymb,amsthm,mathtools}
\usepackage{booktabs,enumitem,xcolor}
\usepackage[numbers,sort&compress]{natbib}
\usepackage[colorlinks=true,linkcolor=blue!55!black,citecolor=blue!55!black,urlcolor=blue!55!black]{hyperref}
\linespread{1.06}
\newtheorem{theorem}{Theorem}[section]
\newtheorem{proposition}[theorem]{Proposition}
\newtheorem{lemma}[theorem]{Lemma}
\newtheorem{corollary}[theorem]{Corollary}
\theoremstyle{remark}\newtheorem{remark}[theorem]{Remark}
\newcommand{\Dtwo}{\operatorname{det}_2}

\title{Resolved Omitted-Shell Control for Regularized Determinants\\
\large A Quantitative Unit-Disk Tail Bound}
\author{Liang Wang}
\date{July 2026}

\begin{document}
\maketitle

\begin{abstract}
The reciprocal local-count audit does not stabilize, but the Arnoldi window
contains additional complete shells beyond every selected RH-222 cloud. This
paper asks whether those resolved omitted shells are small in the
second-regularized determinant.

For $|w|<1$,
\[
 |\log(1-w)+w|\le\frac{|w|^2}{2(1-|w|)}.
\]
Consequently, if a finite omitted resonance multiset $T$ satisfies
$q=R\max_{\lambda\in T}|\lambda|<1$, then on $|z|\le R$,
\[
 \left|\sum_{\lambda\in T}
 [\log(1-z\lambda)+z\lambda]\right|
 \le \frac{R^2}{2(1-q)}
 \sum_{\lambda\in T}|\lambda|^2.
\]
This gives a branch-consistent uniform bound for the omitted
$\det_2$ logarithm.

Each physical endpoint has 12--14 complete resolved roots beyond the selected
cloud. On the closed unit disk, $q\le0.29818$. The maximum analytic upper
bound is $0.16804$, the maximum observed 192-point grid tail is $0.07532$,
and every bound has positive slack.

The result controls the candidate-window tail only. It does not bound the
unresolved infinite operator complement, and therefore does not establish a
uniform small-noise determinant family.
\end{abstract}

\section{Why second regularization helps}

For a resonance multiset $T$, the finite second-regularized factor is
\begin{equation}\label{eq:factor}
 P_{T,2}(z)=
 \prod_{\lambda\in T}(1-z\lambda)e^{z\lambda}.
\end{equation}
The exponential removes the linear term in the logarithm. Near $z=0$, the
tail is therefore quadratic in the omitted resonances rather than linear.

This is the finite-product counterpart of the Hilbert--Schmidt determinant
used at fixed noise in RH-7 \citep{WangRH7,Simon2005}. It is exactly the
regularization required when square summability is available but absolute
summability is not.

\section{The logarithmic tail theorem}

Choose the logarithm by analytic continuation from zero on a disk where no
factor vanishes.

\begin{lemma}[One-factor estimate]\label{lem:one}
For $|w|<1$,
\begin{equation}\label{eq:one}
 |\log(1-w)+w|
 \le\frac{|w|^2}{2(1-|w|)}.
\end{equation}
\end{lemma}
\begin{proof}
The convergent power series gives
\[
 \log(1-w)+w=-\sum_{m=2}^{\infty}\frac{w^m}{m}.
\]
Since $m^{-1}\le1/2$ for $m\ge2$,
\[
 \sum_{m=2}^{\infty}\frac{|w|^m}{m}
 \le\frac12\sum_{m=2}^{\infty}|w|^m
 =\frac{|w|^2}{2(1-|w|)}.
\]
\end{proof}

\begin{theorem}[Uniform finite $\det_2$ tail]\label{thm:tail}
Let $T$ be a finite resonance multiset and let
\[
 q=R\max_{\lambda\in T}|\lambda|<1.
\]
Then the logarithm of \eqref{eq:factor}, normalized to vanish at zero,
satisfies
\begin{equation}\label{eq:bound}
 \sup_{|z|\le R}|\log P_{T,2}(z)|
 \le B_R(T):=
 \frac{R^2}{2(1-q)}\sum_{\lambda\in T}|\lambda|^2.
\end{equation}
\end{theorem}
\begin{proof}
Apply Lemma~\ref{lem:one} to $w=z\lambda$. Since
$|z\lambda|\le q$,
\[
 |\log(1-z\lambda)+z\lambda|
 \le\frac{|z|^2|\lambda|^2}{2(1-q)}.
\]
Sum over $T$ and use $|z|\le R$.
\end{proof}

\begin{corollary}[Multiplicative tail]\label{cor:mult}
On the same disk,
\[
 |P_{T,2}(z)-1|\le e^{B_R(T)}-1.
\]
\end{corollary}
\begin{proof}
Write $P_{T,2}=e^L$ and use
$|e^L-1|\le e^{|L|}-1$.
\end{proof}

The logarithmic form is more informative when composing factors and comparing
channels.

\section{Resolved selected and omitted shells}

At one endpoint, let
\[
 \Lambda_{\rm cand}
 =\Lambda_{\rm sel}\sqcup T_{\rm res}
\]
be the complete-shell part of the Arnoldi candidate window, split after the
RH-222 selected cloud. A nonreal candidate root whose conjugate lies outside
the Arnoldi window is excluded from both sides; no synthetic partner is added.

Then
\begin{equation}\label{eq:split}
 P_{\Lambda_{\rm cand},2}
 =P_{\Lambda_{\rm sel},2}P_{T_{\rm res},2},
\end{equation}
and Theorem~\ref{thm:tail} controls the exact ratio of the two finite products.

Each endpoint contains 12--14 roots in $T_{\rm res}$. The count varies because
real singleton shells and a possible incomplete candidate-boundary root alter
parity.

\section{Physical unit-disk audit}

The disk radius is frozen at $R=1$. The observed logarithmic tail is sampled
at 48 angles on each of the radii
\[
 0.25,\quad0.5,\quad0.75,\quad1.
\]
The analytic bound itself holds at every point of the closed unit disk; the
grid is used only to assess its slack.

\begin{center}
\begin{tabular}{lr}
\toprule
diagnostic & result\\
\midrule
endpoint count & 32\\
resolved omitted roots & 12--14\\
grid points per endpoint & 192\\
maximum $q$ & $0.298172$\\
maximum $\sum_{T_{\rm res}}|\lambda|^2$ & archived\\
maximum logarithmic upper bound & $0.168037$\\
maximum observed grid log tail & $0.075315$\\
minimum bound slack & $0.011053$\\
\bottomrule
\end{tabular}
\end{center}

Every omitted reciprocal zero has modulus at least $q^{-1}>3.35$ when $q$ is
taken at its worst endpoint, so the unit disk is safely zero-free for every
resolved omitted factor.

\section{What the result does not bound}

The full bulk spectrum decomposes as
\[
 \Lambda_{\rm bulk}
 =\Lambda_{\rm sel}\sqcup T_{\rm res}\sqcup T_{\rm unres}.
\]
Theorem~\ref{thm:tail} controls $T_{\rm res}$ because its roots are explicitly
stored. It supplies no value for
\[
 \sum_{\lambda\in T_{\rm unres}}|\lambda|^2.
\]
At fixed noise, Hilbert--Schmidt theory guarantees that the infinite sum is
finite and the truncations converge locally uniformly
\citep{WangRH7,Simon2005}. The missing requirement is a bound uniform as
$\sigma\downarrow0$ after the chosen moving factors have been removed.

A whole-matrix Frobenius norm gives one formal upper bound on the unresolved
eigenvalue square sum. The next paper tests it. Because the operators are
nonnormal, that bound sees singular-direction mass that may be invisible to
the resonance product and can be extremely pessimistic.

\section{Claim boundary}

The exact advance is a quantitative omitted-shell theorem and its successful
application to every resolved candidate window. It does not certify Arnoldi
roots, control the unresolved complement, or prove a small-noise normal
family. The failed local-count gate of RH-227 is not repaired merely by a
small resolved tail \citep{WangRH227}. Gate A remains open.

\bibliographystyle{plainnat}
\bibliography{references}
\end{document}
